\paragraph{Content class of the order-\(\LUV^2\) comparison: the ledger at \(N_s=6\), \(n_D=2\) and the scope of the edge-term question (conditional; not closure).}
Keep the setting of the paragraphs on the normalisation of the tensor kinetic term, on fermionic content, on the order \(\Lambda^2\) and on covariant variables: four matrix coordinates, the Poisson limit taken before the proper-time limit, the Euclidean section, \(\alpha=1\), one loop, the \(\LUV^2\) term with one common proper-time limit, \(N_s\) scalar components, the connection with its ghosts and \(n_D\) Dirac fields with the Yukawa couplings dropped, and the two modes with \(\delta\rho=0\) about a constant \(\bar\Theta\). No state of the matrix model enters: the replicas below are gluing prescriptions for the Poisson-limit operators on the cover of the matrix chart, and no limit of finite matrices or of the rank of a block is taken.

\emph{Sector data.} For each sector \(\sqrt{|G|}\operatorname{tr}(\tfrac16R_G-E)=\tfrac n6\,\rho R[g_\rho]+\kappa\,\rho^2X\) pointwise, with \((n,\kappa)=(1,0)\) for a scalar component and for a ghost, \((4,1)\) for the connection and \((4,\tfrac12)\) for a Dirac field (endomorphisms computed from the operators \(\mathcal O_s\), \(\mathcal O_v\), \(\mathcal O_f\) at four random two-jets of a closed \(\omega\), exact arithmetic), and on the modes with \(\delta\rho=0\) the second-order averages obey \(\langle\rho R[g_\rho]\rangle=-\tfrac14\langle\rho^2X\rangle\) (four random modes, exact arithmetic). With the weights \(1,1,-2,-1\) this gives \(c_R=N_{\mathrm{tot}}/6\), \(\kappa_X=1-\tfrac12n_D\) and \(c_T=c_R-4\kappa_X\). The area coefficients of the preceding paragraphs are \(c_A^{\mathrm{per}}=c_R\) for the unrotated periodic gluing, \(c_A=c_R+n_D\) for the fixed-state replica and \(c_A^{\mathrm{sm}}=-c_R/2\) for the first-order smooth replica, and the Noether coefficient is \(c_W(r)=c_R+\kappa_Xr\). Hence, for every \(N_s\) and \(n_D\),
\[
  c_A-c_T=4\kappa_X+n_D,\qquad c_A^{\mathrm{per}}-c_T=4\kappa_X,\qquad c_W(r)-c_T=\kappa_X(r+4),\qquad (c_A-1)-c_W(-1)=\tfrac12n_D .
\]
The gap of the fixed-state replica has two parts of different origin: the index-weighted coefficient of \(\rho^2X\), which enters the tensor side only, and the excess \(+1\) per Dirac field of the antiperiodic cone value, which enters the area side only. The identification of that gap with \(4\kappa_X\), the relation \(c_A-1=c_W(-1)\) and the remainder \(3\kappa_X\) of the paragraphs on covariant variables and on the matrix gauge realisation hold at \(n_D=0\), as stated there, and at no other \(n_D\).
\begin{center}
\begin{tabular}{@{}lccc@{}}
\toprule
coefficient & general & \((N_s,n_D)=(6,0)\) & \((6,2)\)\\
\midrule
\(c_R=N_{\mathrm{tot}}/6\) & \(\tfrac16(N_s+2-4n_D)\) & \(\tfrac43\) & \(0\)\\
\(\kappa_X\) & \(1-\tfrac12n_D\) & \(1\) & \(0\)\\
\(c_T\) & \(c_R-4\kappa_X\) & \(-\tfrac83\) & \(0\)\\
\(c_T^{\mathrm{ext}}\) & \(c_R-2\kappa_X\) & \(-\tfrac23\) & \(0\)\\
\(c_W(r)\) & \(c_R+\kappa_Xr\) & \(\tfrac43+r\) & \(0\)\\
\(c_A^{\mathrm{per}}\) & \(c_R\) & \(\tfrac43\) & \(0\)\\
\(c_A\), fixed-state & \(c_R+n_D\) & \(\tfrac43\) & \(2\)\\
\(c_A\), fixed-state and vector holonomy, \(n<1\) & \(c_R+n_D-1\) & \(\tfrac13\) & \(1\)\\
\(c_A^{\mathrm{sm}}\) & \(-\tfrac12c_R\) & \(-\tfrac23\) & \(0\)\\
\bottomrule
\end{tabular}
\end{center}

\emph{At \(N_s=6\), \(n_D=2\).} Here \(c_R=\kappa_X=0\), the weighted \(a_1\) density vanishes pointwise for every Poisson bivector, and the induced density of order \(\LUV^2\) is zero. The tensor side then vanishes in every presentation, \(c_T=c_T^{\mathrm{ext}}=c_W(r)=0\) for every \(r\): the family \(L_r\) of the paragraph on covariant variables has a single member, and the dependence of the Noether coefficient on the presentation, the non-invariance of \(\rho^2X\) under local rotations of the frame and the term \(-4\kappa_X\) required there of the modular Hamiltonian are all absent, the remainder \(3\kappa_X\) being zero. The uniform contact \(-N_{\mathrm{tot}}/4\) is zero, and the periodic and smooth replicas give \(0=0\). The one nonvanishing entry is the fixed-state area coefficient \(c_A=n_D=2\), which is the Dirac excess alone; the vector holonomy contact of modulus \(1\), the edge-mode value of~\cite{DonnellyWall2016} for a Maxwell field if the state carries that sector, changes it to \(1\) on \(n<1\) and \(3\) on \(n>1\), and the Dirac holonomy contact vanishes at first order, so no contact of the paragraph on the order \(\Lambda^2\) removes it. The coefficient of \(W\) at logarithmic order has the same pattern, with gaps \(0\), \(2\), \(0\) for the periodic, fixed-state and completed replicas and \(\kappa_C^T=0\); and since \(\mathcal S_4\) is quartic in \(b_\sigma\), the logarithmic coefficient of the induced action stated in the paragraph on the next order has no quadratic form about a constant \(\bar\Theta\). At this content the requirement \(G_A=G_T\) of the paragraph on the normalisation is therefore, at order \(\LUV^2\), the single condition \(c_A=0\) on the area side, with \(1/G_T=0\): for the trace of \(\rho^n\) of the paragraph on fermion replica gluing \(1/G_A=\LUV^2/\pi\) while \(1/G_T=0\), and for the periodic and smooth replicas both vanish. In this class the comparison at this content is either an equality of vanishing coefficients or the inequality \(2\neq0\), and in neither case does it determine a Newton constant.

\emph{Sectors with a scale.} Let each sector carry a constant mass, \(\mathcal O_i\to\mathcal O_i+m_i^2\), the formal shift of the paragraph on the next order, with the ghosts at the mass of the connection. The heat trace of the sector is multiplied by \(e^{-sm_i^2}\) on the chart and on the cover alike, so the sector data above are weighted on both sides by the same function, \(f(m^2)=\int_{\LUV^{-2}}^\infty\mathrm ds\,s^{-2}e^{-sm^2}=\LUV^2-m^2\ln(\LUV^2/m^2)-(1-\gamma_{\mathrm E})m^2+O(m^4/\LUV^2)\), with \(\gamma_{\mathrm E}\) the Euler constant. With \(\mathsf S_n=\sum_iw_in_if(m_i^2)\) and \(\mathsf S_\kappa=\sum_iw_i\kappa_if(m_i^2)\), in units in which the massless value of \(f\) is \(1\),
\[
  c_T=\tfrac16\mathsf S_n-4\mathsf S_\kappa,\qquad c_A^{\mathrm{per}}=\tfrac16\mathsf S_n,\qquad c_A=\tfrac16\mathsf S_n+\sum_{\text{Dirac}}f(m_D^2),\qquad c_A^{\mathrm{per}}-c_T=4\mathsf S_\kappa .
\]
At \(N_s=6\), \(n_D=2\) the \(\LUV^2\) parts of \(\mathsf S_n\) and \(\mathsf S_\kappa\) cancel and both sums are of order \(m^2\ln\LUV\). The tensor coefficient is nonzero only if \(\mathsf S_n\neq24\,\mathsf S_\kappa\), and for the periodic gluing it equals the area coefficient if and only if \(\mathsf S_\kappa=0\), a condition on the connection and Dirac spectra alone; the fixed-state replica keeps in addition the Dirac term, of order \(\LUV^2\). For one bosonic mass and one fermionic mass \(c_T=-\tfrac83(f_b-f_f)\) and \(c_A^{\mathrm{per}}=\tfrac43(f_b-f_f)\), of opposite sign. At this content the comparison of \(G_A\) with \(G_T\) is thus not vacuous only in a class with a split spectrum, where it reduces to the two sums \(\mathsf S_n\) and \(\mathsf S_\kappa\). Whether the Kaluza--Klein towers of fuzzy extra dimensions, for which~\cite{Steinacker2023} attributes the Einstein--Hilbert term to a mixed quartic term and not to a mass shift of \(a_1\), give sums of this form is not examined.

These statements are for the Poisson limit taken first, four matrix coordinates, the Euclidean section, one loop, \(\alpha=1\), the proper-time regulator with one common limit, the contents stated with the Yukawa couplings dropped, the replicas and contacts of the preceding paragraphs with their hypotheses, the two modes with \(\delta\rho=0\), and formal constant masses without a gauge-invariant mass term. They supply no state, modular Hamiltonian or gluing of the matrix model, do not derive the content, do not treat finite terms, Moyal corrections, more than four matrix coordinates or a Kaluza--Klein tower, use the stated properties of \(\mathcal S_4\) and of the coefficient of \(W\) from the preceding paragraphs without a new check, and give no value or sign of \(G_{\mathrm{eff}}\). Problem~B remains open and PROGRAMME-CONDITIONAL.

